\section{Linear models in business, science, and engineering}
A model begins by defining its variables and their units. Matrix columns
then describe what one unit of each input contributes, and rows record
one output requirement at a time. A mathematically consistent nutritional
system is not physically feasible if it requires negative amounts of food.
In a circuit, signs depend on the chosen loop directions and battery
polarities. In a migration model, column $j$ records where people starting
in region $j$ go, so its entries sum to $1$ when nobody leaves the system.
The labels [M] in the text permit matrix software; exact reduced matrices
and substitution checks are supplied below rather than treating a software
answer as a substitute for interpreting the model.

\practice{1.10.P1}{Express the diet example as a matrix equation}
Let $x_1,x_2,x_3$ be the amounts of nonfat milk, soy flour, and whey,
respectively, measured in $100$-gram units. Put the nutrient contributions
per unit into columns and the targets into the right-hand side:
\[
 A=\begin{bmatrix}36&51&13\\52&34&74\\0&7&1.1\end{bmatrix},\qquad
 x=\begin{bmatrix}x_1\\x_2\\x_3\end{bmatrix},\qquad
 b=\begin{bmatrix}33\\45\\3\end{bmatrix},\qquad Ax=b.
\]
The rows specify protein, carbohydrate, and fat in grams. For example,
the first row says $36x_1+51x_2+13x_3=33$, so it adds the protein
from all three ingredients. The actual food masses in grams are $100x_j$.

\exercise{1.10.1}{A two-cereal mixture with four requirements}
\step{(a) Define the serving counts and the vector equation.}
Let $x$ be servings of Cheerios and $y$ servings of Natural Cereal.
With rows ordered as calories, protein, carbohydrate, and fat,
\[
 x\begin{bmatrix}110\\4\\20\\2\end{bmatrix}
 +y\begin{bmatrix}130\\3\\18\\5\end{bmatrix}
 =\begin{bmatrix}295\\9\\48\\8\end{bmatrix}.
\]
\step{(b) Solve and check all the requirements.}
The equivalent matrix equation is
\[
 \begin{bmatrix}110&130\\4&3\\20&18\\2&5\end{bmatrix}
 \begin{bmatrix}x\\y\end{bmatrix}
 =\begin{bmatrix}295\\9\\48\\8\end{bmatrix}.
\]
Use the protein equation $4x+3y=9$ and fat equation $2x+5y=8$.
Subtract twice the latter from the former: $-7y=-7$, hence $y=1$.
Then $2x+5=8$ gives $x=3/2$. Check the other rows:
$110(3/2)+130=295$ and $20(3/2)+18=48$. Both serving
counts are nonnegative, so the solution is physically meaningful.
\answer{Mix $1.5$ servings of Cheerios with $1$ serving of Natural Cereal.}

\exercise{1.10.2}{A specified mixture and a new nutritional target}
\step{(a) Build the contribution matrix.}
Order the rows as calories, protein, fiber, and fat, and the columns
as Shredded Wheat and Crispix. Then
\[
 B=\begin{bmatrix}160&110\\5&2\\6&0.1\\1&0.4\end{bmatrix},\quad
 u=\begin{bmatrix}3\\2\end{bmatrix},\quad
 Bu=\begin{bmatrix}700\\19\\18.2\\3.8\end{bmatrix}.
\]
For instance the fiber amount is $6(3)+0.1(2)=18.2$ grams.
\step{(b) Solve for the requested different mixture.}
The desired equation is $B(x,y)^T=(130,3.20,2.46,0.64)^T$.
Its reduced augmented matrix is
\[
 \begin{bmatrix}1&0&0.4\\0&1&0.6\\0&0&0\\0&0&0\end{bmatrix}.
\]
Thus $x=0.4$ and $y=0.6$. To check all rows, the totals are
$64+66=130$ calories, $2+1.2=3.20$ protein, $2.4+0.06=2.46$
fiber, and $0.4+0.24=0.64$ fat. These nonnegative proportions
also add to one serving.
\answer{The target is attainable with $0.4$ serving of Shredded Wheat
and $0.6$ serving of Crispix.}

\exercise{1.10.3}{Compare two versions of a lunch}
Let $x,y,z$ be servings of the pasta dish, broccoli, and chicken.
Use calories, protein, and fiber as the three rows. Solving each system
requires all three targets, not only the calories.
\step{(a) Use the original Mac and Cheese.}
\[
 \begin{bmatrix}270&51&70\\10&5.4&15\\2&5.2&0\end{bmatrix}
 \begin{bmatrix}x\\y\\z\end{bmatrix}
 =\begin{bmatrix}400\\30\\10\end{bmatrix}.
\]
Exact row reduction, treating the decimals as the stated rational
numbers, gives
\[
 [A\mid b]\sim
 \begin{bmatrix}1&0&0&8205/8323\\0&1&0&12850/8323\\0&0&1&6550/8323\end{bmatrix}.
\]
Hence $(x,y,z)\approx(0.98582,1.54391,0.78698)$. Every variable
is determined by a pivot, so the solution is unique. Substituting the
exact fractions into the original three rows gives $(400,30,10)^T$.
All three fractions are positive, so the mixture is feasible.
\step{(b) Replace the pasta column by the whole-wheat-shells data.}
Only that ingredient's column changes:
\[
 \begin{bmatrix}260&51&70\\9&5.4&15\\5&5.2&0\end{bmatrix}
 \begin{bmatrix}x\\y\\z\end{bmatrix}
 =\begin{bmatrix}400\\30\\10\end{bmatrix}.
\]
Reduction gives $(x,y,z)=(5470,4400,5180)/5023$, approximately
$(1.08899,0.87597,1.03126)$. Substitution again gives exactly
$400$ calories, $30$ grams protein, and $10$ grams fiber before
rounding. The broccoli portion drops from about $1.54$ to $0.88$
serving, in line with the stated preference.
\answer{To two decimal places the mixtures are $(0.99,1.54,0.79)$
and $(1.09,0.88,1.03)$ servings in the stated ingredient order.
Use the exact fractions for exact nutritional totals.}

\exercise{1.10.4}{A consistent diet equation can be physically impossible}
\step{(a) Specify the four variables and four nutrient rows.}
Let $x_1,x_2,x_3,x_4$ be $100$-gram units of nonfat milk, soy flour,
whey, and isolated soy protein. In the row order protein,
carbohydrate, fat, calcium, the equation is
\[
 \begin{bmatrix}
 36&51&13&80\\52&34&74&0\\0&7&1.1&3.4\\1.26&0.19&0.8&0.18
 \end{bmatrix}
 \begin{bmatrix}x_1\\x_2\\x_3\\x_4\end{bmatrix}
 =\begin{bmatrix}33\\45\\3\\0.8\end{bmatrix}.
\]
Each coefficient is grams of that nutrient per $100$ grams of ingredient.
\step{(b) Solve and interpret the signs.}
Exact reduction yields the unique solution
\[
 x=\begin{bmatrix}
 2264687/3531104\\1921341/3531104\\-326887/3531104\\-1468529/7062208
 \end{bmatrix}
 \approx\begin{bmatrix}0.64135\\0.54412\\-0.09257\\-0.20794\end{bmatrix}.
\]
Multiplication of the exact vector by the original matrix reproduces
all four targets, so the algebraic equations are consistent. But the
whey and isolated-protein amounts are negative. One cannot subtract
food already absent from the mixture, so this unique algebraic solution
is not a permissible recipe. Because it is unique, there is no different
nonnegative solution to these same equations.
\answer{These four ingredients cannot meet all four targets with
nonnegative amounts under the supplied nutrient data.}

\textbf{Circuit convention for Exercises 5--8.}
Use the counterclockwise loop directions drawn in the textbook.
For each loop, add all resistances on its boundary to get the diagonal
coefficient. A shared resistor carries the difference of the adjacent
loop currents, so its resistance enters the neighboring column with a
minus sign. On the right, crossing a battery from its short plate to
its long plate contributes a positive voltage; the reverse crossing
contributes a negative voltage. These rules explain every sign in
$Ri=v$. A negative current means the actual direction is opposite
to the assumed counterclockwise direction.


\exercise{1.10.5}{Derive the loop equations and solve them}

The four diagonal resistance sums are $1+1+4+5=11$, $3+5+1+1=10$, $4+1+2+2=9$, and $1+2+3+4=10$. The successive shared resistances are $5,1,2$. Battery crossings give $20+30=50$, $-30-10=-40$, $10+20=30$, and $-20-10=-30$ volts.

Thus, in amperes, the matrix equation is
\[ \begin{bmatrix}11 & -5 & 0 & 0\\-5 & 10 & -1 & 0\\0 & -1 & 9 & -2\\0 & 0 & -2 & 10\end{bmatrix}\begin{bmatrix}I_{1}\\I_{2}\\I_{3}\\I_{4}\end{bmatrix}=\begin{bmatrix}50\\-40\\30\\-30\end{bmatrix}.\]

Exact elimination of this augmented system gives a pivot in every variable column and hence the unique solution
\[ i=\begin{bmatrix}\frac{265}{72}\\- \frac{137}{72}\\\frac{185}{72}\\- \frac{179}{72}\end{bmatrix}\approx \begin{bmatrix}3.6806\\-1.9028\\2.5694\\-2.4861\end{bmatrix}.\]

The displayed exact fractions satisfy $Ri=v$ in every row. For example, the first row computes $50$, the first specified net voltage. All remaining rows were checked by the same substitution. Interpret a minus sign as a direction reversal rather than a negative physical magnitude.


\exercise{1.10.6}{Derive the loop equations and solve them}

The diagonal sums are $3+2+1=6$, $1+4+4=9$, $4+1+2=7$, and $2+2+3=7$. Adjacent loops share resistances $1,4,2$. Every battery on the left is crossed from short plate to long plate in the chosen direction, giving voltages $30,20,40,10$.

Thus, in amperes, the matrix equation is
\[ \begin{bmatrix}6 & -1 & 0 & 0\\-1 & 9 & -4 & 0\\0 & -4 & 7 & -2\\0 & 0 & -2 & 7\end{bmatrix}\begin{bmatrix}I_{1}\\I_{2}\\I_{3}\\I_{4}\end{bmatrix}=\begin{bmatrix}30\\20\\40\\10\end{bmatrix}.\]

Exact elimination of this augmented system gives a pivot in every variable column and hence the unique solution
\[ i=\begin{bmatrix}\frac{3630}{571}\\\frac{4650}{571}\\\frac{6700}{571}\\\frac{2730}{571}\end{bmatrix}\approx \begin{bmatrix}6.3573\\8.1436\\11.734\\4.7811\end{bmatrix}.\]

The displayed exact fractions satisfy $Ri=v$ in every row. For example, the first row computes $30$, the first specified net voltage. All remaining rows were checked by the same substitution. Interpret a minus sign as a direction reversal rather than a negative physical magnitude.


\exercise{1.10.7}{Derive the loop equations and solve them}

Keep the order $I_1$ upper left, $I_2$ lower left, $I_3$ lower right, $I_4$ upper right. The diagonal sums are $1+7+4=12$, $7+6+2=15$, $6+5+3=14$, and $4+5+4=13$. Shared resistances are $7$ between loops 1 and 2, $6$ between 2 and 3, $5$ between 3 and 4, and $4$ between 4 and 1. The source voltages are $40,30,20,-10$; the last is negative because loop 4 crosses that battery against its positive direction.

Thus, in amperes, the matrix equation is
\[ \begin{bmatrix}12 & -7 & 0 & -4\\-7 & 15 & -6 & 0\\0 & -6 & 14 & -5\\-4 & 0 & -5 & 13\end{bmatrix}\begin{bmatrix}I_{1}\\I_{2}\\I_{3}\\I_{4}\end{bmatrix}=\begin{bmatrix}40\\30\\20\\-10\end{bmatrix}.\]

Exact elimination of this augmented system gives a pivot in every variable column and hence the unique solution
\[ i=\begin{bmatrix}\frac{119910}{10487}\\\frac{110640}{10487}\\\frac{84270}{10487}\\\frac{61240}{10487}\end{bmatrix}\approx \begin{bmatrix}11.434\\10.55\\8.0357\\5.8396\end{bmatrix}.\]

The displayed exact fractions satisfy $Ri=v$ in every row. For example, the first row computes $40$, the first specified net voltage. All remaining rows were checked by the same substitution. Interpret a minus sign as a direction reversal rather than a negative physical magnitude.


\exercise{1.10.8}{Derive the loop equations and solve them}

The loop order is the four outer regions $1,2,3,4$, then the central diamond $5$. Their diagonal sums are $3+1+4+1=9$, $1+1+3+2=7$, $2+2+3+3=10$, $1+1+2+3=7$, and $4+3+3+2=12$. The central loop shares resistances $4,3,3,2$ with the four outer loops. Additional shared stems are $1$ between 1 and 4, $1$ between 1 and 2, $2$ between 2 and 3, and $3$ between 3 and 4. Its own battery contribution is zero; the outer contributions are $50,-30,20,-40$.

Thus, in amperes, the matrix equation is
\[ \begin{bmatrix}9 & -1 & 0 & -1 & -4\\-1 & 7 & -2 & 0 & -3\\0 & -2 & 10 & -3 & -3\\-1 & 0 & -3 & 7 & -2\\-4 & -3 & -3 & -2 & 12\end{bmatrix}\begin{bmatrix}I_{1}\\I_{2}\\I_{3}\\I_{4}\\I_{5}\end{bmatrix}=\begin{bmatrix}50\\-30\\20\\-40\\0\end{bmatrix}.\]

Exact elimination of this augmented system gives a pivot in every variable column and hence the unique solution
\[ i=\begin{bmatrix}\frac{3019}{755}\\- \frac{3309}{755}\\- \frac{683}{755}\\- \frac{4382}{755}\\- \frac{722}{755}\end{bmatrix}\approx \begin{bmatrix}3.9987\\-4.3828\\-0.90464\\-5.804\\-0.95629\end{bmatrix}.\]

The displayed exact fractions satisfy $Ri=v$ in every row. For example, the first row computes $50$, the first specified net voltage. All remaining rows were checked by the same substitution. Interpret a minus sign as a direction reversal rather than a negative physical magnitude.


\textbf{Migration convention for Exercises 9--13.}
Put the current populations into a column vector. In each migration
matrix, column $j$ tells what fractions of starting population $j$
arrive in the various destinations. Thus $x_{k+1}=Mx_k$ and two
years later means $x_2=M(Mx_0)$, not $Mx_0$.

\exercise{1.10.9}{A seven-percent outward migration rate}
Let $x_k=(c_k,s_k)^T$ denote city and suburban populations $k$
years after 2010. The matrix and initial population are
\[
 M=\begin{bmatrix}0.93&0.05\\0.07&0.95\end{bmatrix},\qquad
 x_0=\begin{bmatrix}800000\\500000\end{bmatrix},\qquad x_{k+1}=Mx_k.
\]
For example $c_{k+1}=0.93c_k+0.05s_k$: city residents who stay
plus suburban residents who move into the city. First multiplication
gives $x_1=(769000,531000)^T$. Multiplying once more gives
\[
 x_2=\begin{bmatrix}741720\\558280\end{bmatrix}.
\]
The total remains $1300000$, as a check on the column sums and arithmetic.
\answer{In 2012 the model predicts $741720$ city residents and
$558280$ suburban residents.}

\exercise{1.10.10}{A six-percent outward migration rate}
With city first and suburbs second,
\[
 M=\begin{bmatrix}0.94&0.04\\0.06&0.96\end{bmatrix},\quad
 x_0=\begin{bmatrix}10000000\\800000\end{bmatrix},\quad x_{k+1}=Mx_k.
\]
Here the first column says $94\%$ of city residents stay and
$6\%$ move to the suburbs. The second says $4\%$ move in and
$96\%$ stay in the suburbs. Successive multiplications give
\[
 x_1=\begin{bmatrix}9432000\\1368000\end{bmatrix},\qquad
 x_2=\begin{bmatrix}8920800\\1879200\end{bmatrix}.
\]
Both totals equal the initial $10800000$.
\answer{The 2012 prediction is $8920800$ in the city and $1879200$
in the suburbs, subject to the model's exclusions.}

\exercise{1.10.11}{Estimate migration rates before iterating}
\step{(a) Divide each flow by its population of origin.}
The fraction leaving California is
$516100/31524000\approx0.0163716533$, hence $0.01637$ to five
decimal places. The fraction moving into California from elsewhere
is $381262/228680000\approx0.0016672293$, hence $0.00167$.
Put California first. The rounded migration model is
\[
 M=\begin{bmatrix}0.98363&0.00167\\0.01637&0.99833\end{bmatrix},\qquad
 x_0=\begin{bmatrix}31524000\\228680000\end{bmatrix}.
\]
The diagonal entries are the complements of the corresponding
outward rates, not independently rounded unrelated quantities.
\step{(b) Apply six annual steps to reach 2000.}
Compute $x_{k+1}=Mx_k$ six times, retaining full arithmetic with
the specified five-place coefficients. The result is
\[
 x_6\approx\begin{bmatrix}30754526.4814\\229449473.5186\end{bmatrix}.
\]
To the nearest person, these are $30754526$ in California and
$229449474$ elsewhere. Their sum is the unchanged total $260204000$.
The fractional entries before rounding are population estimates, not
literal fractional people. Using unrounded migration rates would yield
slightly different estimates, so it is important to state which rates were used.
\answer{The model predicts about $30.755$ million in California and
$229.449$ million elsewhere, with births, deaths, and international
migration excluded as required.}

\exercise{1.10.12}{Follow a fleet through two daily transitions}
Use the order airport, east, west. The columns correspond to the
rental origin and rows to the return location, so
\[
 M=\begin{bmatrix}0.97&0.05&0.10\\0&0.90&0.05\\0.03&0.05&0.85\end{bmatrix},
 \qquad x_0=\begin{bmatrix}295\\55\\150\end{bmatrix}.
\]
Monday to Tuesday gives $x_1=Mx_0=(303.9,57,139.1)^T$.
Tuesday to Wednesday gives
\[
 x_2=Mx_1=\begin{bmatrix}311.543\\58.255\\130.202\end{bmatrix}.
\]
For example the first entry is $0.97(303.9)+0.05(57)+0.10(139.1)$.
The three entries sum to $500$ before rounding, matching the fleet size.
\answer{Approximately $312$ cars at the airport, $58$ at the east
location, and $130$ at the west location on Wednesday.}

\exercise{1.10.13}{Twenty iterations and their limiting pattern}
The source model uses
$M=\begin{bmatrix}0.95&0.03\\0.05&0.97\end{bmatrix}$.
The total population is $1000000$ in both parts. Therefore, if $c_k$
is the city population, the suburb population is $1000000-c_k$,
and the first row of the recurrence becomes
\[
 c_{k+1}=0.95c_k+0.03(1000000-c_k)=0.92c_k+30000.
\]
Subtract $375000$ from both sides to obtain
$c_{k+1}-375000=0.92(c_k-375000)$. Repeated substitution gives
\[
 c_k=375000+(c_0-375000)(0.92)^k,\qquad s_k=1000000-c_k.
\]
This formula both checks the matrix iterations and explains their pattern.
\step{(a) Start from $(600000,400000)^T$.}
The city component is $375000+225000(0.92)^k$.
\step{(b) Start from $(350000,650000)^T$.}
The city component is $375000-25000(0.92)^k$.
The following table lists all twenty requested vectors, with entries
rounded to the nearest person only for display.
\begin{center}
\begin{tabular}{r|rr|rr}
 &\multicolumn{2}{c|}{Part (a)}&\multicolumn{2}{c}{Part (b)}\\
$k$&City&Suburbs&City&Suburbs\\\hline

1&582000&418000&352000&648000\\
2&565440&434560&353840&646160\\
3&550205&449795&355533&644467\\
4&536188&463812&357090&642910\\
5&523293&476707&358523&641477\\
6&511430&488570&359841&640159\\
7&500515&499485&361054&638946\\
8&490474&509526&362170&637830\\
9&481236&518764&363196&636804\\
10&472737&527263&364140&635860\\
11&464918&535082&365009&634991\\
12&457725&542275&365808&634192\\
13&451107&548893&366544&633456\\
14&445018&554982&367220&632780\\
15&439417&560583&367843&632157\\
16&434264&565736&368415&631585\\
17&429522&570478&368942&631058\\
18&425161&574839&369427&630573\\
19&421148&578852&369872&630128\\
20&417456&582544&370283&629717\\

\end{tabular}\end{center}
Since $0<0.92<1$, its powers tend to zero. In (a) the city
population decreases toward $375000$, passing below the suburban
population between years $7$ and $8$ (first below at the year-$8$ update).
In (b) it increases toward the same
$375000$. The suburb populations approach $625000$ from the
opposite sides. The common limiting proportions come from the
migration rule, not from the initial distribution.
\answer{Both initial distributions tend to $(375000,625000)^T$,
but they approach it from different sides.}

\exercise{1.10.14}{Temperature models obey superposition}
\step{(a) Write one averaging equation for each interior point.}
Order the points as $1$ upper left, $2$ upper right, $3$ lower right,
$4$ lower left, as in the diagrams. Moving the two interior-neighbor
temperatures to the left gives the common coefficient matrix
\[
 K=\begin{bmatrix}4&-1&0&-1\\-1&4&-1&0\\0&-1&4&-1\\-1&0&-1&4\end{bmatrix}.
\]
For plate A, each point touches boundary temperatures $20$ and $0$,
so $KT_A=(20,20,20,20)^T$. Direct substitution gives
$T_A=(10,10,10,10)^T$.
For plate B, the two relevant boundary values sum to $10,40,50,20$
in point order. Thus $KT_B=(10,40,50,20)^T$, and solving gives
\[
 T_B=(10,17.5,20,12.5)^T.
\]
For example its first equation is $4(10)-17.5-12.5=10$,
and its third is $4(20)-17.5-12.5=50$.
Their sum is
\[
 T_A+T_B=(20,27.5,30,22.5)^T,
\]
exactly the temperature list cited from Section 1.1. The boundary
data there are the sums of these two plates' boundary data, so the
same additivity appears at the interior points.
\step{(b) Triple all of plate A's boundary temperatures.}
The new right-hand side is $(60,60,60,60)^T$, three times the
old one. Hence $T=(30,30,30,30)^T=3T_A$. This is also
checked directly: each row is $4(30)-30-30=60$.
\step{(c) Formulate and justify the general rule.}
The map from the eight boundary values to the four interior values
is linear: adding boundary data adds interior solutions, and scaling
boundary data scales the interior solution. To justify uniqueness,
if two solutions existed, their difference $z$ would satisfy $Kz=0$.
Choose a component whose absolute value is largest, say $|z_i|=L$.
Its equation gives $4L\le2L$, since there are two neighboring
components, each of absolute value at most $L$. Thus $L=0$ and
$z=0$. There cannot be two distinct solutions.
Existence for every right-hand side $b$ can be verified directly by
substitution in the explicit solution formula
\[
 T=\frac1{24}
 \begin{bmatrix}7&2&1&2\\2&7&2&1\\1&2&7&2\\2&1&2&7\end{bmatrix}b.
\]
Each $b_i$ is a sum of two boundary values, so both stages of the
construction use only fixed linear combinations.
\answer{The interior-temperature map is linear. Its additivity
explains part (a), and its homogeneity explains part (b).}
